\section{Repaso del Diffusion Model: de la imagen a la biología}

\subsection{Principio central del modelo de difusión}

Antes de profundizar en el diseño de proteínas, Amini primero repasó la idea central del Diffusion Model. El principio básico del modelo de difusión es \textbf{lograr la generación aprendiendo a recuperar los datos a partir del ruido}, e incluye dos procesos inversos entre sí:

\textbf{1. Proceso directo de adición de ruido (Forward Noising Process)}

Partiendo del espacio de datos original (como una imagen), se destruyen los datos de manera gradual mediante la adición iterativa de ruido:

$$x_0 \xrightarrow{\text{noise}} x_1 \xrightarrow{\text{noise}} x_2 \xrightarrow{\text{noise}} \cdots \xrightarrow{\text{noise}} x_T \sim \mathcal{N}(0, I)$$

Donde:
\begin{itemize}
    \item $x_0$: datos originales limpios
    \item $x_t$: datos después de agregar ruido en el paso $t$
    \item $x_T$: después de $T$ pasos, se aproxima a ruido puramente aleatorio
    \item Este proceso \textbf{no requiere entrenamiento}; simplemente se agrega ruido gaussiano de manera gradual
\end{itemize}

\textbf{2. Proceso inverso de eliminación de ruido (Reverse Denoising Process)}

Se entrena una red neuronal para aprender el proceso inverso, recuperando los datos de manera gradual a partir del ruido:

$$x_T \xrightarrow{\text{denoise}} x_{T-1} \xrightarrow{\text{denoise}} \cdots \xrightarrow{\text{denoise}} x_0$$

La tarea de entrenamiento puede formalizarse así: dado el dato con ruido $x_t$ en el paso de tiempo $t$, predecir una versión con menos ruido $x_{t-1}$ o, de manera equivalente, predecir el ruido residual $\epsilon$:

$$\mathcal{L} = \mathbb{E}_{x_0, \epsilon, t} \left[ \| \epsilon - \epsilon_\theta(x_t, t) \|^2 \right]$$

Donde:
\begin{itemize}
    \item $\epsilon$: el ruido realmente agregado
    \item $\epsilon_\theta$: el ruido predicho por la red neuronal
    \item $t$: el paso de tiempo muestreado aleatoriamente
\end{itemize}

\begin{importantbox}{¿Por qué el Diffusion Model es adecuado para el diseño de proteínas?}
El Diffusion Model ha mostrado tres ventajas revolucionarias en el campo de la generación de imágenes:
\begin{enumerate}
    \item \textbf{Calidad de generación}: el realismo y la fidelidad de las muestras son extremadamente altos
    \item \textbf{Diversidad}: puede cubrir una amplia región de la distribución de datos, sin sufrir mode collapse
    \item \textbf{Controlabilidad}: admite generación condicional, lo que permite especificar el objetivo de generación
\end{enumerate}
Estas características corresponden precisamente a las necesidades centrales del diseño de proteínas: las proteínas generadas deben ser biológicamente viables (calidad), cubrir funciones diversas (diversidad) y satisfacer especificaciones de diseño particulares (controlabilidad).
\end{importantbox}

\subsection{El reto de pasar del dominio continuo al discreto}

El Diffusion Model estándar procesa \textbf{datos continuos} (como los valores de los píxeles de una imagen), y el proceso directo se implementa agregando ruido gaussiano continuo. Sin embargo, la secuencia de proteínas es \textbf{datos discretos} —en cada posición hay uno de 20 aminoácidos posibles— y no se le puede simplemente «agregar ruido gaussiano» a un carácter.

Esto plantea un reto técnico central: \textbf{¿cómo definir el proceso de difusión sobre datos discretos?}

\subsection{Resumen de la sección}

El núcleo del Diffusion Model consiste en aprender el proceso inverso que va del ruido a los datos. Su éxito en el ámbito de las imágenes ofrece un paradigma sólido del cual tomar referencia para el diseño de proteínas, pero la naturaleza discreta de las secuencias de proteínas exige una transformación fundamental del marco estándar de difusión continua.

